\documentclass[11pt]{article}
\usepackage[a4paper,margin=27mm]{geometry}
\usepackage{amsmath,amssymb,amsthm,hyperref}
\newtheorem{proposition}{Proposition}
\newtheorem{lemma}[proposition]{Lemma}
\newcommand{\E}{\mathbb E}
\newcommand{\Pp}{\mathbb P}
\newcommand{\Pois}{\operatorname{Pois}}
\newcommand{\1}{\mathbf1}
\title{Independent review of the two-jump Poisson comparison}
\author{}
\date{28 October 2026}
\begin{document}
\maketitle

This note proves component estimates for a direct count comparison. It does
not prove a new range of the Berry--Esseen inequality. It identifies a false
allocation shortcut in the earlier discovery program, gives the correct
finite-count variables, and establishes a finite evaluation formula and a
useful monotonicity. The accepted small-noise theorems and the local
variance-aware moment lemma do not use the false shortcut.

\section{Comparison with two Poisson jump sizes}
Write $Z_b(V)=b(P_{V/b^2}-V/b^2)$, where $P_a\sim\Pois(a)$.
Different occurrences are independent unless stated otherwise; $Z_b(0)=0$.
We use the following published theorem: if independent $X_j$ satisfy
$\E X_j\le0$, $X_j\le b$, and $\sum_j\E X_j^2\le V$, then
\[
 \E f\Bigl(\sum_jX_j\Bigr)\le\E f(Z_b(V))
 \quad\text{for } f\in\mathcal F^2_+.
\]
Here $\mathcal F^2_+$ consists of continuously differentiable functions
whose function and derivative are increasing and convex. This is
Pinelis, \emph{On the Bennett--Hoeffding inequality}, Proposition~2.8
and Corollary~2.11, printed p.~9,
\href{https://arxiv.org/pdf/0902.4058}{arXiv:0902.4058}.
In particular $f(x)=(x-y)_+^2$ belongs to this class. The premise is a
variance upper bound, not an equality. Zero variance is dealt with directly.

\begin{proposition}[Groupwise comparison]
Suppose $S_0,S_1$ are sums of two independent groups of centered independent
variables. The summands in group $j$ are at most $b_j>0$, and their total
variance is at most $V_j$. Then, for every real $y$,
\begin{equation}\label{eq:groups}
 \E(S_0+S_1-y)_+^2
 \le \E\{Z_{b_0}(V_0)+Z_{b_1}(V_1)-y\}_+^2.
\end{equation}
Consequently, for every $y<b$,
\begin{equation}\label{eq:closed}
 \Pp(S_0+S_1\ge b)
 \le \frac{\E\{Z_{b_0}(V_0)+Z_{b_1}(V_1)-y\}_+^2}{(b-y)^2}.
\end{equation}
\end{proposition}
\begin{proof}
Condition on $S_1=z$ and apply the published theorem to the translated
hinge $(x+z-y)_+^2$. Integrate this inequality. Then condition on the
new independent variable $Z_{b_0}(V_0)$ and replace the second group.
All integrands are nonnegative, so Tonelli applies; all expectations
are finite by the second-moment bounds. No theorem for a larger class
of moment functions is being assumed. Finally
$\1_{[b,\infty)}(x)\le (x-y)_+^2/(b-y)^2$. This includes an atom at
$b$ and requires the strict inequality $y<b$.
\end{proof}
For the incoming lower tail of $W_j$, apply this with $S=-W_j$:
the jump bound is $A$ on $-U_0$ and $d\ge a_1$ on $-U_1$.
The inward bound on $U_0$ is not a permissible replacement for $A$.

\section{The allocation belongs to the marginal, not to each count}
Let $p\in(0,1)$, $q=1-p$, and let $U_0,U_1$ be centered conditional
laws with variances $v_0,v_1$. Assume
$-A\le U_0\le a_0$, $-a_1\le U_1\le A$, with
$a_0,a_1\ge0$ and $a_0+a_1\le d$. Set
\[
 s=qv_0+pv_1>0,\quad \lambda=ns,\quad
 w=\frac{pv_1}{s}\in[0,1],\quad k=np+\delta\in\{0,\ldots,n\}.
\]
At a legal count $k+i$, the exact group variances are
\begin{align}\label{eq:transport}
 T_{0,i}&=\lambda(1-w)\left(1-\frac{\delta+i}{nq}\right),&
 T_{1,i}&=\lambda w\left(1+\frac{\delta+i}{np}\right).
\end{align}
These follow simply from $v_0=s(1-w)/q$, $v_1=sw/p$ and the counts
$n-k-i,k+i$. Thus $t_i=\lambda+(\delta+i)(v_1-v_0)$, but generally
$T_{0,i}\ne(1-w)t_i$ and $T_{1,i}\ne wt_i$.

If $p,q\ge2/5$ and $|\delta|\le D$, then the common inflation
\[
 T_{0,i}\le\lambda(1-w)\left(1+\frac{5(D+i)}{2n}\right),\qquad
 T_{1,i}\le\lambda w\left(1+\frac{5(D+i)}{2n}\right)\quad(i\ge1)
\]
is valid. The earlier discovery program's incoming inflation therefore
has a legitimate interpretation. Its own-count/cubic-repair join did not:
there it replaced the two own-count variances by $(1-w)t,wt$ while using
$w$ as the marginal allocation in the moment repair.

\paragraph{An exact admissible counterexample to that intermediate bound.}
Take $n=42$, $p=1/2$, $k=22$, $r=2/25$, $A=3/5$, and
$\varepsilon=1/1000$. Let $U_0$ be the centered two-point law on
$\{-A,\varepsilon\}$, with probabilities
$\varepsilon/(A+\varepsilon),A/(A+\varepsilon)$, and let $U_1=-U_0$.
This gives an actual four-point marginal $B+U$, not a moment relaxation.
Here
\[
 v_0=v_1=v=3/5000,\quad w=1/2,\quad \lambda=63/2500,
 \quad T_0=20v,\quad T_1=22v,\quad t=42v.
\]
With $\kappa=-4r$, put $c_0=A^2-4Ar$, $c_1=A^2+4Ar$.
The valid bound is
$\E W^4-\kappa\E W^3\le c_0T_0+c_1T_1+3t^2$.
The common-fraction substitute fails by the exact amount
\begin{equation}\label{eq:false}
 \E W^4-\kappa\E W^3-
 \{[c_0(1-w)+c_1w]t+3t^2\}
 =\frac{423903}{2500000000}>0.
\end{equation}
Indeed $\E U_0^3=v(\varepsilon-A)$,
$\E U_0^4=v(A^2-A\varepsilon+\varepsilon^2)$, and independence gives
$\E W^4=42\E U_0^4+3t^2-126v^2$.
The difference is
$v[8rA-8r\varepsilon-168A\varepsilon+42\varepsilon^2]$.
The inward sum is $.002<.072$. The raw threshold at phase $u=-r$ is
\[
 |H|=\frac{\delta-r}{\sqrt{pq+s}}
      =\frac{23/25}{\sqrt{.2506}}<1.84.
\]
Thus the counterexample satisfies the intended support and threshold
constraints. It invalidates an intermediate discovery inference, not the
sharp conjecture, the accepted $.03$ theorem, or the local moment lemma.

\paragraph{The corrected joint cubic prices.}
For a quartic minorant
$h(x)=(-x^4+b_3x^3+b_2x^2+b_1x+b_0)/M$, $M>0$, define
$\kappa=b_3-4r$, $L=b_2+3b_3r-6r^2$, $J=Mh(r)$.
Put $\Lambda_0=\lambda(1-w)$, $\Lambda_1=\lambda w$,
$g_0=1-\delta/(nq)$, $g_1=1+\delta/(np)$, so
$t=g_0\Lambda_0+g_1\Lambda_1$.
Let $m>0$, $C\ge0$, and assume $a_0<p$, $a_1<q$. Write
$\mu_j=\E U_j^3$ and $\ell=(pq+s)^{-1/2}$. These support bounds fix
the signs of $B+U-p$, giving
\[
 \rho=\ell^3\{pq(1-2pq)+3pq(v_0+v_1)+p\mu_1-q\mu_0\}.
\]
Set $v=pq$ and $\rho_B=(1-2v)/\sqrt v$. Choose
\[
 B\ge \frac{v(1-2v)}s\{(1+s/v)^{3/2}-1\},
\]
and suppose the resulting lower bound for the entire cubic bracket is
nonnegative.
If $0\le K\le C\ell^3/n$, one may lower-bound that entire repair
using $K$. For the right-panel bounds, the bracket is positive since
$3\min(p,q)-B-d>0$. Combining the third moments before estimating them
then gives the lower bound
\begin{align*}
 m\Pp(W>-r)+C(\rho-\rho_B)\ \ge\ &
 \frac mM(Lt-3t^2+J)+K\lambda\{3p(1-w)+3qw-B\}\\
 &-\Lambda_0\max_{x\in\{-A,a_0\}}
       \left\{\frac mM g_0(x^2-\kappa x)+Kx\right\}\\
 &-\Lambda_1\max_{x\in\{-a_1,A\}}
       \left\{\frac mM g_1(x^2-\kappa x)-Kx\right\}.
\end{align*}
To see this, integrate each displayed quadratic against
$x^2\,dP_j(x)/v_j$ and use its endpoint maximum; it is convex because
$g_j\ge0$. At $v_j=0$ the corresponding term is zero directly.
This preserves the two physical multipliers. Replacing $K$ in separated
signed pieces without first proving the whole bracket nonnegative is
not justified. For fixed other parameters the displayed bound is concave
in $w$, owing to $-3t^2$; the endpoint prices are independent of $w$.

\section{Monotonicity removes some allocation work, but not all}
\begin{proposition}[Variance and allocation monotonicity]\label{prop:mono}
For $0<d\le A$, $T\ge0$, $w\in[0,1]$, and $y\in\mathbb R$, set
\[
 G(T,w;y)=\E\{Z_A(T(1-w))+Z_d(Tw)-y\}_+^2.
\]
Then $G$ is nondecreasing in $T$ and nonincreasing in $w$.
\end{proposition}
\begin{proof}
First prove $\E(Z_d(V)-y)_+^2\le\E(Z_A(V)-y)_+^2$.
For $V>0$ and integers $N\ge V/d^2$, take independent
$B_j\sim\operatorname{Bernoulli}(V/(Nd^2))$.
The centered sum $X_N=d\sum_{j=1}^N(B_j-\E B_j)$ has summands bounded
above by $d\le A$ and total variance at most $V$. The published comparison
gives the required bound with $X_N$ in place of $Z_d(V)$.
The binomial Poisson limit gives $X_N\Rightarrow Z_d(V)$.
Moreover
\[
 \E X_N^4\le d^2V+3V^2,
\]
since each centered summand has absolute value at most $d$ and independence
gives the fourth-moment expansion. Hence the quadratic hinges are uniformly
integrable: their tails are controlled by a constant times
$(1+|X_N|)^2$, whose squares have uniformly bounded expectation.
Truncating a continuous hinge and then removing the truncation justifies
the limit. At $V=0$ the assertion is immediate.

The comparison persists after adding an independent variable with finite
second moment: condition on that variable and integrate the translated
hinge inequalities. For $w_1<w_2$, use a common sum consisting of
$Z_A(T(1-w_2))$ and $Z_d(Tw_1)$. The remaining variance is
$V=T(w_2-w_1)$, assigned to large jumps at $w_1$ and small jumps at
$w_2$. The preceding comparison proves monotonicity in $w$.
For $T_1<T_2$, the law at $T_2$ is the law at $T_1$ plus independent
centered Poisson increments with variances
$(T_2-T_1)(1-w)$ and $(T_2-T_1)w$. Conditional Jensen proves
monotonicity in $T$. All assertions include zero intensities.
\end{proof}

For example, choose valid upper group variances
$V_0=\lambda(1-w)\omega_0$, $V_1=\lambda w\omega_1$ with
$0\le\omega_0\le\omega_1$. Their small-jump fraction is at least $w$.
If $V_0+V_1\le T_+$ and $w\ge w_-$, the incoming hinge is bounded by
$G(T_+,w_-;y)$. One must check $\omega_0\le\omega_1$; it does not follow
from arbitrary count multipliers without a phase condition.
On a negative phase, $\delta=\xi+r$, $|\xi|<1$, so $\delta+i>0$
for $i\ge1$ and one may take $\omega_0=1$ and a suitable $\omega_1\ge1$.

\paragraph{Convexity is false, even for a positive hinge anchor.}
At $A=3/5$, $d=9/125$, $T=1/40$, $w=1/4$, and $y=1/5$,
\begin{equation}\label{eq:curv}
 \frac{\partial^2}{\partial w^2}G(T,w;y)
 =-e^{-815/648}\,
 \frac{787398879114798488065}{147382613147155571933184}<0.
\end{equation}
Here is an exact finite verification. Put
$\mu=T(1-w)/A^2$, $\nu=Tw/d^2$, $M=A\mu+d\nu$, $z=M+y$.
At the stated point $\mu+\nu=815/648$, $z=229/720$, and
$4d<z<5d<A$. In a neighborhood, the only atoms below the hinge are
zero large jumps and $j=0,\ldots,4$ small jumps. Thus
\[
 G(T,w;y)=T+y^2-e^{-\mu-\nu}Q(w),\qquad
 Q(w)=\sum_{j=0}^4\frac{\nu^j}{j!}(M+y-jd)^2.
\]
Since $(\mu+\nu)'=385/81$ and $M'=11/36$, differentiating this finite
polynomial shows that the coefficient in \eqref{eq:curv} is
$Q''-2(385/81)Q'+(385/81)^2Q$. The independent rational checker
reconstructs it. Therefore a concave own-count bound minus this hinge
need not be concave. Allocation endpoints alone cannot be inferred from
convexity. Monotone bounds on closed allocation intervals remain valid.

\section{The complete hinge is a finite complement}
For $A=3/5$, $d=9/125$, set $a=3/125$, so $A=25a$, $d=3a$.
Let $S=25P_\mu+3P_\nu$. For every real $K$,
\begin{equation}\label{eq:lattice}
 \E(S-K)_+^2=625\mu+9\nu+(25\mu+3\nu-K)^2
 -e^{-\mu-\nu}\sum_{25j+3l\le K}
 \frac{\mu^j\nu^l}{j!l!}(K-25j-3l)^2.
\end{equation}
The sum is finite (empty if $K<0$). This is the identity
$x_+^2=x^2-(-x)_+^2$, followed by the exact first two moments of $S$.
It includes all upper-tail mass and every lattice collision; there is no
truncation error. The physical anchor is
$K=(y+A\mu+d\nu)/a$, and the physical numerator is $a^2$ times
\eqref{eq:lattice}. A rational $K$ is allowed.

The independent implementation uses the probability-generating function
$e^{\mu z^{25}+\nu z^3}=\sum_{j\ge0}c_jz^j$ instead of a double sum:
\[
 c_0=1,\qquad j c_j=25\mu c_{j-25}+3\nu c_{j-3},
 \quad c_j=0\ (j<0).
\]
The finite deficit is $\sum_{j\le K}(K-j)^2c_j$.
For rational $x=\mu+\nu\ge0$, let $S_N=\sum_{j=0}^Nx^j/j!$ and
$N+2>x$. The positive exponential series gives
\[
 S_N\le e^x\le S_N+\frac{x^{N+1}}{(N+1)!}\frac1{1-x/(N+2)}.
\]
Reciprocals enclose $e^{-x}$. A lower exponential bound gives an upper
hinge numerator because the deficit is subtracted. This direction is
essential. It remains valid at zero intensity and at integer $K$.

For an alternative box bound with fixed $K$, the uncentered numerator
$\E(S-K)_+^2$ increases with each intensity by ordinary Poisson coupling.
Use upper intensity endpoints in that numerator and a positive lower
bound on $b/a+25\mu+3\nu-K$ in the denominator. A centered physical
anchor generally moves $K$ with the intensities; in that case one should
apply Proposition~\ref{prop:mono}, or account for this movement explicitly.

\section{Bounding all remaining incoming counts}
Every legal count has variance at most
$n\max(v_0,v_1)\le(5/2)\lambda$ when $p,q\ge2/5$.
Put $\mu=(5/2)\lambda/A^2$ and assume $0<A<1$, $r\ge0$.
For an integer $i\ge2$, the Poisson hinge at integer anchor $i$ gives
\[
 \Pp(W_{k+i}\le-i-r)
 \le\frac{\mu^{i+1}}{(i+1)!D_i^2},\qquad
 D_i=i(1/A-1)+r/A+\mu>0.
\]
Indeed $(j-i)_+^2\le\binom j{i+1}$ for integers $j\ge0$ and $i\ge2$.
For $j=i+l$, $l\ge1$, the ratio
$\binom{i+l}{i+1}/l^2$ starts at one and is nondecreasing, since
\[
 \frac{\binom{i+l+1}{i+1}/(l+1)^2}
      {\binom{i+l}{i+1}/l^2}
 =\frac{l(i+l+1)}{(l+1)^2}\ge1
 \quad\Longleftrightarrow\quad (i-1)l\ge1.
\]
Poisson factorial moments give the numerator. Since $D_i$ increases,
for any $m\ge2$ with $\mu<m+2$ the complete tail is bounded by
\[
 \sum_{i=m}^{n-k}\Pp(W_{k+i}\le-i-r)
 \le\frac{\mu^{m+1}}{(m+1)!D_m^2}\frac1{1-\mu/(m+2)}.
\]
The added terms beyond legal counts are nonnegative numbers in the
majorant; no nonexistent conditional law is invoked. The condition
$i\ge2$ matters: at $i=1,j=3$, the proposed pointwise bound would read
$4\le3$. At $\lambda=0$ handle the deterministic conditional noises
directly.

\paragraph{Status.}
These are complete component derivations under the stated hypotheses.
The published Poisson comparison remains an external theorem. Exact checks
cover the transport counterexample, a genuine failure of convexity,
finite-complement arithmetic, actual conditional convolutions and rejected
shortcuts. They are not a closed-domain count cover, nor an independent
proof of the published theorem. No old certificate or manuscript file is
changed.
\end{document}
